%%
%% This is file `sample-sigconf.tex',
%% generated with the docstrip utility.
%%
%% The original source files were:
%%
%% samples.dtx  (with options: `all,proceedings,bibtex,sigconf')
%% 
%% IMPORTANT NOTICE:
%% 
%% For the copyright see the source file.
%% 
%% Any modified versions of this file must be renamed
%% with new filenames distinct from sample-sigconf.tex.
%% 
%% For distribution of the original source see the terms
%% for copying and modification in the file samples.dtx.
%% 
%% This generated file may be distributed as long as the
%% original source files, as listed above, are part of the
%% same distribution. (The sources need not necessarily be
%% in the same archive or directory.)
%%
%%
%% Commands for TeXCount
%TC:macro \cite [option:text,text]
%TC:macro \citep [option:text,text]
%TC:macro \citet [option:text,text]
%TC:envir table 0 1
%TC:envir table* 0 1
%TC:envir tabular [ignore] word
%TC:envir displaymath 0 word
%TC:envir math 0 word
%TC:envir comment 0 0
%%
%% The first command in your LaTeX source must be the \documentclass
%% command.
%%
%% For submission and review of your manuscript please change the
%% command to \documentclass[manuscript, screen, review]{acmart}.
%%
%% When submitting camera ready or to TAPS, please change the command
%% to \documentclass[sigconf]{acmart} or whichever template is required
%% for your publication.
%%
%%

\documentclass[sigconf, anonymous]{acmart}
\usepackage{xspace}
\usepackage{algorithm}
\usepackage[noend]{algorithmic}
\usepackage{amsmath}
\usepackage{amsthm}

\newcommand{\sys}{CARAMEL\xspace}
\newcommand{\alg}{\textsc{CaramelSort}\xspace}

\newcommand{\KL}[2]{D_{\mathrm{KL}}\!\left(#1 \;\middle\|\; #2\right)}
\newcommand{\Hb}[1]{H\!\left(#1\right)}
\newcommand{\pbar}{\bar{p}}
\newcommand{\bits}{\text{ bits}}
\newcommand{\Lmax}{L_{\max}}
\newcommand{\bL}{b_L}

%%
%% \BibTeX command to typeset BibTeX logo in the docs
\AtBeginDocument{%
  \providecommand\BibTeX{{%
    Bib\TeX}}}

%% Rights management information.  This information is sent to you
%% when you complete the rights form.  These commands have SAMPLE
%% values in them; it is your responsibility as an author to replace
%% the commands and values with those provided to you when you
%% complete the rights form.
\setcopyright{acmlicensed}
\copyrightyear{2018}
\acmYear{2018}
\acmDOI{XXXXXXX.XXXXXXX}
%% These commands are for a PROCEEDINGS abstract or paper.
\acmConference[Conference acronym 'XX]{Make sure to enter the correct
  conference title from your rights confirmation email}{June 03--05,
  2018}{Woodstock, NY}
%%
%%  Uncomment \acmBooktitle if the title of the proceedings is different
%%  from ``Proceedings of ...''!
%%
%%\acmBooktitle{Woodstock '18: ACM Symposium on Neural Gaze Detection,
%%  June 03--05, 2018, Woodstock, NY}
\acmISBN{978-1-4503-XXXX-X/2018/06}


%%
%% Submission ID.
%% Use this when submitting an article to a sponsored event. You'll
%% receive a unique submission ID from the organizers
%% of the event, and this ID should be used as the parameter to this command.
%%\acmSubmissionID{123-A56-BU3}

%%
%% For managing citations, it is recommended to use bibliography
%% files in BibTeX format.
%%
%% You can then either use BibTeX with the ACM-Reference-Format style,
%% or BibLaTeX with the acmnumeric or acmauthoryear sytles, that include
%% support for advanced citation of software artefact from the
%% biblatex-software package, also separately available on CTAN.
%%
%% Look at the sample-*-biblatex.tex files for templates showcasing
%% the biblatex styles.
%%

%%
%% The majority of ACM publications use numbered citations and
%% references.  The command \citestyle{authoryear} switches to the
%% "author year" style.
%%
%% If you are preparing content for an event
%% sponsored by ACM SIGGRAPH, you must use the "author year" style of
%% citations and references.
%% Uncommenting
%% the next command will enable that style.
%%\citestyle{acmauthoryear}


%%
%% end of the preamble, start of the body of the document source.
\begin{document}

%%
%% The "title" command has an optional parameter,
%% allowing the author to define a "short title" to be used in page headers.
\title{\sys: Sweetening Compressed Static Functions for Space-Efficient Caching in Search and Recommendation Systems}

% \shorttitle{\sys: Space-Efficient Caching in Search and Recommendation Systems}

%%
%% The "author" command and its associated commands are used to define
%% the authors and their affiliations.
%% Of note is the shared affiliation of the first two authors, and the
%% "authornote" and "authornotemark" commands
%% used to denote shared contribution to the research.
\author{Ben Trovato}
\authornote{Both authors contributed equally to this research.}
\email{trovato@corporation.com}
\orcid{1234-5678-9012}
\author{G.K.M. Tobin}
\authornotemark[1]
\email{webmaster@marysville-ohio.com}
\affiliation{%
  \institution{Institute for Clarity in Documentation}
  \city{Dublin}
  \state{Ohio}
  \country{USA}
}

\author{Lars Th{\o}rv{\"a}ld}
\affiliation{%
  \institution{The Th{\o}rv{\"a}ld Group}
  \city{Hekla}
  \country{Iceland}}
\email{larst@affiliation.org}

\author{Valerie B\'eranger}
\affiliation{%
  \institution{Inria Paris-Rocquencourt}
  \city{Rocquencourt}
  \country{France}
}

\author{Aparna Patel}
\affiliation{%
 \institution{Rajiv Gandhi University}
 \city{Doimukh}
 \state{Arunachal Pradesh}
 \country{India}}

\author{Huifen Chan}
\affiliation{%
  \institution{Tsinghua University}
  \city{Haidian Qu}
  \state{Beijing Shi}
  \country{China}}

\author{Charles Palmer}
\affiliation{%
  \institution{Palmer Research Laboratories}
  \city{San Antonio}
  \state{Texas}
  \country{USA}}
\email{cpalmer@prl.com}

\author{John Smith}
\affiliation{%
  \institution{The Th{\o}rv{\"a}ld Group}
  \city{Hekla}
  \country{Iceland}}
\email{jsmith@affiliation.org}

\author{Julius P. Kumquat}
\affiliation{%
  \institution{The Kumquat Consortium}
  \city{New York}
  \country{USA}}
\email{jpkumquat@consortium.net}

%%
%% By default, the full list of authors will be used in the page
%% headers. Often, this list is too long, and will overlap
%% other information printed in the page headers. This command allows
%% the author to define a more concise list
%% of authors' names for this purpose.
\renewcommand{\shortauthors}{Trovato et al.}
\renewcommand{\shorttitle}{Space-Efficient Result Caching for Search and Recommendation Systems}
%%
%% The abstract is a short summary of the work to be presented in the
%% article.
\begin{abstract}
We consider the problem of space-efficient result caching in search and recommender systems. A \emph{result cache} is a key-value data store mapping entities like queries or user ids to a list of associated items. With the heavy computational cost of modern retrieval pipelines, particularly in the age of large language models (LLMs), precomputing and caching results is often a necessity. However, result caching shifts the resource burden from compute to storage. Prior work on result caching has not addressed the distinct challenge that such caches must store a \emph{list} of values per key, rather than a single entity.
Our key insight is that result caches often exhibit strong power-law behavior, with frequent repetition of items across different keys. We propose to exploit this property by modeling the cache as a matrix — where rows correspond to keys and columns to ranked positions — and building \emph{compressed static functions} for each column. 

This architecture raises open research questions about codebook sharing across columns, entropy-reducing permutations for unordered results, and the handling of ragged lists of varying length.
In this paper, we systematically analyze each of these questions through a theoretical and experimental lens, culminating in the design of \sys, a new system for low-memory result caching in retrieval pipelines. By extending compressed static functions to handle lists of values, \sys performs lookups directly on highly compressed representations of the result cache while maintaining microsecond-level query latency. On workloads derived from real production data of a major e-commerce platform, \sys achieves an XXX reduction in memory compared to in-memory caching baselines, yielding a Pareto-optimal tradeoff between latency and memory footprint across the systems evaluated.
\end{abstract}

%%
%% The code below is generated by the tool at http://dl.acm.org/ccs.cfm.
%% Please copy and paste the code instead of the example below.
%%
\begin{CCSXML}
<ccs2012>
 <concept>
  <concept_id>00000000.0000000.0000000</concept_id>
  <concept_desc>Do Not Use This Code, Generate the Correct Terms for Your Paper</concept_desc>
  <concept_significance>500</concept_significance>
 </concept>
 <concept>
  <concept_id>00000000.00000000.00000000</concept_id>
  <concept_desc>Do Not Use This Code, Generate the Correct Terms for Your Paper</concept_desc>
  <concept_significance>300</concept_significance>
 </concept>
 <concept>
  <concept_id>00000000.00000000.00000000</concept_id>
  <concept_desc>Do Not Use This Code, Generate the Correct Terms for Your Paper</concept_desc>
  <concept_significance>100</concept_significance>
 </concept>
 <concept>
  <concept_id>00000000.00000000.00000000</concept_id>
  <concept_desc>Do Not Use This Code, Generate the Correct Terms for Your Paper</concept_desc>
  <concept_significance>100</concept_significance>
 </concept>
</ccs2012>
\end{CCSXML}

\ccsdesc[500]{Do Not Use This Code~Generate the Correct Terms for Your Paper}
\ccsdesc[300]{Do Not Use This Code~Generate the Correct Terms for Your Paper}
\ccsdesc{Do Not Use This Code~Generate the Correct Terms for Your Paper}
\ccsdesc[100]{Do Not Use This Code~Generate the Correct Terms for Your Paper}

%%
%% Keywords. The author(s) should pick words that accurately describe
%% the work being presented. Separate the keywords with commas.
\keywords{Do, Not, Use, This, Code, Put, the, Correct, Terms, for,
  Your, Paper}
%% A "teaser" image appears between the author and affiliation
%% information and the body of the document, and typically spans the
%% page.

\received{20 February 2007}
\received[revised]{12 March 2009}
\received[accepted]{5 June 2009}

%%
%% This command processes the author and affiliation and title
%% information and builds the first part of the formatted document.
\maketitle

\section{Introduction}

Modern search and recommendation pipelines face a fundamental tension. Low online query latency is a critical pillar of user experience, yet retrieval algorithms are growing increasingly computationally expensive — a trend accelerated by the recent adoption of large language models (LLMs). To reconcile these competing demands, search and recommendation systems turn to result caching. A result cache is a key-value lookup table mapping entities such as search queries or user IDs to a list of associated retrieved items. By precomputing and caching results, a retrieval system can bypass real-time query processing and model inference, which can be prohibitively slow and expensive. Below, we highlight three motivating examples of result caching in practice.

\noindent \textbf{Example 1:} Homepage Recommendations. A user logs in to a video streaming site and sees movie recommendations on the landing page. To generate these suggestions, the recommender engine precomputes a list of movies per user using offline models. Upon login, the system looks up the the list of precomputed results from the cache for that user, optionally reranks this list based on real-time session features, and then serves the top $k$ results to the homepage.   


\noindent \textbf{Example 2:} Item-to-Item Recommendations. When an online shopper views a particular product, an e-commerce service recommends similar or complementary products to drive further engagement. To generate these recommendations, the service precomputes a result cache where the keys are product IDs and the values are the lists of associated product IDs. 

\noindent \textbf{Example 3:} Personalized Search. A search engine uses a personalization model to vary the results served to different users for the same query. Here the keys of the result cache will be \emph{tuples} of queries and user IDs while the values will be lists of search results, (represented, for example, as document IDs).

A common structure emerges across these examples: each key maps to a list of items, naturally forming a matrix. Throughout this paper, we will model result caches as a matrix where the rows correspond to individual users, items, or search queries and the columns represent the list of precomputed results. For a large web service with 500M rows and a precomputed list of 100 items per row (represented as 32 bit IDs), this cache, assuming no further optimizations, will take up at least 200GB of space. In practice, this space footprint might expand even further for cases where the retrieval services maintain multiple such caches, such as one for each shopping category in an e-commerce service. This heavy storage cost motivates the key research question we address in this paper: \textbf{Can we design a space-efficient index for result caching that reduces the storage footprint while retaining fast lookups?} 

We affirmatively answer this question via a new caching system we call \sys (\textbf{C}ompressed \textbf{A}rray \textbf{R}epresentations \textbf{A}nd \textbf{M}ore \textbf{E}fficient \textbf{L}ookups). Our key observation guiding the design of \sys is that there are many common elements across the value lists in result caches. This heavy-tailed popularity phenomenon is well-documented in practical search and recommender systems; some web pages, movies, products, and songs are overwhelmingly more popular than others and thus appear frequently across different keys in the cache. 

To take advantage of this redundancy in the value lists, we leverage \emph{compressed static functions} (CSFs), a succinct retrieval data structure that achieves a space footprint proportional to the empirical entropy of the values with $O(1)$ query time. Viewing result caches as an $n \times m$ matrix, we propose to construct a compressed static function for each \emph{column} of the matrix. Due to the fact that each value list exhibits power law behavior, this architecture can take advantage of repetition within columns to reduce space at the cost of increasing the query time to $O(m)$. In practice, $m$ is typically a small constant (often between 10-100) so we hypothesize that the query latency increase will be negligible, especially when compared to running the full retrieval pipeline in real time.

This idea of building a CSF for each column of the result cache matrix raises a number of open research questions. First, CSFs, which rely on Huffman coding for compression, maintain a codebook to map Huffman codes back to the original item. When we have multiple CSFs, one for each column, is it more space-efficient to share one codebook across all CSFs or maintain separate codebooks per column (\textbf{RQ1})? Secondly, in many cases, the precomputed entries in a result cache are re-ranked at query time to take advantage of real-time personalization signals. Thus, the order in which the results are stored in the cache may not matter. In this order-agnostic setting, a natural optimization is to reorder the entries in each row of the matrix to minimize the sum of columnwise entropies to reduce space even further. However, to our knowledge, the computational complexity of this precise optimization problem has not been studied and it is unclear if there exists an efficient algorithm for this task (\textbf{RQ2}). Finally, a third open question raised by this multi-CSF architecture is how to handle the case of \emph{ragged} value lists where the list of retrieved results per row varies (\textbf{RQ3}).

In this paper, we systematically study each of these research questions through a theoretical and experimental lens, culminating in the design of the \sys system. Our specific contributions are as follows:

\begin{itemize}
\item Problem Statement

\item Codebook sharing

\item Row Permutation

\item Ragged Arrays
\end{itemize}
\section{Background}

A \emph{compressed static function} (CSF)~\cite{hreinsson2009storing} is a data structure
that stores a mapping $f : X \to \Sigma$ from a finite key set $X$ of size $n$ to an
alphabet $\Sigma$, using space per key close to the empirical entropy of the output
values. The core construction encodes each output value with a variable-length
prefix-free codeword and builds a bit array $A$ such that XOR-ing $k$ positions of
$A$, at positions determined by hash functions of the key, recovers the codeword.
This reduction from a function-storage problem to a system of linear equations over
$\mathbb{Z}/2\mathbb{Z}$ (specifically a $k$-XORSAT instance) yields a bit array of
size approximately
\[
    \delta_k \sum_{x \in X} |c(f(x))| \text{bits},
\]
where $c(\cdot)$ is the chosen prefix-free code. 

When the code matches the empirical distribution of outputs, the average codeword
length equals the empirical entropy $H_0$, so the bit array consumes $\approx \delta_k
n H_0 \text{bits}$ in total. Beyond the bit array, one must also store the \emph{codebook}:
the Huffman decoding table and the mapping from codewords to elements of $\Sigma$.


\section{Methodology}

Describe how the system works
\section{Codebook Sharing}

In this section, we study the decision between maintaining separate codebooks for each column in \sys or using a single shared codebook\footnote{In principle, one could also consider a hybrid strategy where a subset of columns share a codebook and the remaining maintain individual codebooks. However, this optimization problem is NP-Hard (via a reduction from \textsc{Set Cover} \cite{karp1972}), so we restrict our attention to the binary decision of full sharing versus no sharing.}. We will proceed by developing cost models for the space usage of each respective index and analyze under what conditions is one setting superior to the other. 

A length-limited canonical Huffman codebook over alphabet $\Sigma$ with maximum
codeword length $\Lmax$ requires storing two things:
\begin{enumerate}
    \item \textbf{Codeword lengths.} One integer in $\{1, \ldots, \Lmax\}$ per symbol,
          each requiring $\bL = \lceil\log_2 \Lmax\rceil$ bits: total
          $|\Sigma|\,\bL$ bits.
    \item \textbf{Symbol ordering.} The canonical ordering maps each codeword to a
          symbol; storing this permutation of $\Sigma$ requires
          $|\Sigma|\lceil\log_2|\Sigma|\rceil$ bits.
\end{enumerate}
The total codebook cost is therefore
\[
    \mathcal{C} = |\Sigma|\bigl(\bL + \lceil\log_2|\Sigma|\rceil\bigr) \bits.
\]
Crucially, since $\Lmax$ is a fixed design constant, $\mathcal{C}$ is the \emph{same}
for both a private and a shared codebook: both store one table over the same alphabet
$\Sigma$ with the same maximum depth $\Lmax$. There is no additional depth penalty from
pooling distributions.

\subsection{All-private strategy}

The total space under strategy $\mathbf{P}$ is
\[
    \mathrm{Space}(\mathbf{P})
    = \underbrace{\delta_k n \sum_{j=1}^{m} H(p_j)}_{\text{bit arrays}}
    + \underbrace{m \cdot \mathcal{C}}_{\text{codebooks}}
    + \underbrace{m \cdot M}_{\text{chunk metadata}},
\]
where $M = (n / 2^\ell) \cdot 64\bits$ is the per-column chunk metadata cost
(one 64-bit word per chunk of approximately $2^\ell$ keys), and
\[
    \mathcal{C} = |\Sigma|\bigl(\bL + \lceil\log_2|\Sigma|\rceil\bigr).
\]

\subsection{All-shared strategy}

By the source-coding theorem, encoding distribution $p_j$ with a code optimised for
$\pbar$ incurs an expected codeword length of
\[
    \ell(p_j, \pbar) = H(p_j) + \KL{p_j}{\pbar}.
\]
Therefore the total space under strategy $\mathbf{S}$ is
\[
    \mathrm{Space}(\mathbf{S})
    = \underbrace{\delta_k n \sum_{j=1}^{m}\bigl[H(p_j) + \KL{p_j}{\pbar}\bigr]}_{\text{bit arrays}}
    + \underbrace{\mathcal{C}}_{\text{one codebook}}
    + \underbrace{m \cdot M}_{\text{chunk metadata}}.
\]

The sum of KL divergences in the shared strategy admits a clean closed form.

\begin{lemma}\label{lem:kl-identity}
    \[
        \sum_{j=1}^{m} \KL{p_j}{\pbar} = m \cdot \Hb{\pbar} - \sum_{j=1}^{m} \Hb{p_j}.
    \]
\end{lemma}

\begin{proof}
By definition,
\[
    \KL{p_j}{\pbar}
    = \sum_{\sigma} p_j(\sigma)\log\frac{p_j(\sigma)}{\pbar(\sigma)}
    = -\Hb{p_j} - \sum_{\sigma} p_j(\sigma)\log\pbar(\sigma).
\]
Summing over $j$ and noting that $\frac{1}{m}\sum_j p_j(\sigma) = \pbar(\sigma)$,
\[
    \sum_{j=1}^{m} \KL{p_j}{\pbar}
    = -\sum_{j=1}^{m}\Hb{p_j}
      - m\sum_{\sigma}\pbar(\sigma)\log\pbar(\sigma)
    = m\Hb{\pbar} - \sum_{j=1}^{m}\Hb{p_j}. \qedhere
\]
\end{proof}

Applying Lemma~\ref{lem:kl-identity}, the bit-array cost under $\mathbf{S}$ simplifies
to
\[
    \delta_k n \sum_{j=1}^{m}\bigl[H(p_j) + \KL{p_j}{\pbar}\bigr]
    = \delta_k n m \Hb{\pbar}.
\]

\begin{remark}
    The total bit-array cost of the shared strategy depends only on the entropy of the
    \emph{pooled} distribution $\pbar$, not on the individual column entropies. Adding a
    high-entropy outlier column to the shared pool is therefore always expensive.
\end{remark}

The two space expressions can now be written compactly as:
\begin{align}
    \mathrm{Space}(\mathbf{P}) &= \delta_k n \sum_{j=1}^{m} \Hb{p_j}
        + m \cdot |\Sigma|\bigl(\bL + \lceil\log_2|\Sigma|\rceil\bigr)
        + m \cdot M, \label{eq:space-P} \\
    \mathrm{Space}(\mathbf{S}) &= \delta_k n m \cdot \Hb{\pbar}
        + |\Sigma|\bigl(\bL + \lceil\log_2|\Sigma|\rceil\bigr)
        + m \cdot M. \label{eq:space-S}
\end{align}

The chunk-metadata terms $m \cdot M$ cancel in the comparison
$\mathrm{Space}(\mathbf{S}) < \mathrm{Space}(\mathbf{P})$. Subtracting
\eqref{eq:space-P} from \eqref{eq:space-S} and applying Lemma~\ref{lem:kl-identity}
gives the following exact, parameter-free criterion.

\begin{proposition}[Decision criterion]\label{prop:criterion}
    Strategy $\mathbf{S}$ (all-shared) uses strictly less space than strategy
    $\mathbf{P}$ (all-private) if and only if
    \[
        \boxed{
            \delta_k n \sum_{j=1}^{m}\KL{p_j}{\pbar}
            \;<\;
            (m-1)\cdot|\Sigma|\bigl(\bL + \lceil\log_2|\Sigma|\rceil\bigr),
        }
    \]
    where $\bL = \lceil\log_2 \Lmax\rceil$ and $\Lmax$ is the fixed maximum codeword
    length.
\end{proposition}

The left-hand side is the \emph{total KL penalty} paid in the bit arrays by using a
mismatched codebook. The right-hand side is the \emph{net codebook saving} from paying
$\mathcal{C}$ once instead of $m$ times; it is a compile-time constant depending only
on $|\Sigma|$ and $\Lmax$.
\section{The \alg Algorithm}


Algorithm~\ref{alg:col-entropy} permutes elements within each row of $A$ to
minimize the total column entropy $\sum_{j=1}^{m} H(A_{:,j})$. Since each
column is compressed independently with its own codebook, a column's storage
cost is bounded by its Shannon entropy; minimizing the sum therefore directly
minimizes aggregate storage. The algorithm proceeds in two phases: a fast
frequency-based initialization followed by iterative hill-climbing refinement.

\textbf{Phase~1} sorts each row by descending global frequency. Because
high-frequency values appear in many rows, this causes them to concentrate in
the same leftmost columns across all rows, immediately producing low-entropy
columns. The sort is independent per row and trivially parallelizable, running
in $O(Nm\log m)$ total.

\textbf{Phase~2} iteratively improves column assignments. Let $c_v^{(j)}$
denote the number of rows that place value $v$ in column $j$. The entropy of
column $j$ decomposes as
\begin{equation}\label{eq:entropy-decomp}
  H(A_{:,j}) \;=\; \log N \,-\, \frac{1}{N}\sum_v c_v^{(j)} \log c_v^{(j)}\,,
\end{equation}
so minimizing total column entropy is equivalent to maximizing the global
concentration
\begin{equation}\label{eq:objective}
  F(\sigma) \;=\; \sum_{j=1}^{m} \sum_v c_v^{(j)} \log c_v^{(j)}\,.
\end{equation}
Because $g(x) = x\log x$ is strictly convex, its marginal $g(c+1) - g(c)$
increases in $c$: moving a copy of value $v$ from a column where it appears
$b$ times to one where it already appears $a > b$ times strictly increases
$F$. The locally optimal greedy rule is therefore to place each item in the
column where its count is already largest.

This motivates the refinement strategy. At each iteration, a
\emph{scoreboard} of the current counts $c_v^{(j)}$ is built, and each value's
columns are ranked by count into a preference list $\mathrm{pref}(v)$. Every
row is then reassigned independently via a greedy max-heap: items compete for
their top-preferred columns and fall back to the next preference on conflict.
The per-row step is $O(m\log m)$ and embarrassingly parallel; the scoreboard
is rebuilt from scratch in $O(Nm)$ per iteration, avoiding stale counts.

The algorithm converges rapidly in practice---typically within 10--15
iterations---because the frequency sort already provides a near-optimal
initialization, leaving the refinement phase to resolve only local conflicts.
We detect convergence by a relative entropy change threshold $\varepsilon$
(e.g., $0.1\%$) and terminate early.




\begin{algorithm}[t]
\caption{CARAMELSort Heuristic (Draft)}
\label{alg:col-entropy}
\begin{algorithmic}
  \STATE {\bfseries Input:} $A \in \mathbb{Z}^{N \times m}$ without duplicates in each row, max iterations $T$, threshold $\varepsilon$
  \STATE {\bfseries Output:} Permuted $A \in \mathbb{Z}^{N \times m}$ minimizing $\sum_{j=1}^{m} H(A_{:,j})$
  \STATE {\bfseries Phase 1: Global Frequency Sort}
  \STATE $\mathrm{freq}(v) = |\{(i,j) : A_{i,j} = v\}|$ for each value $v$
  \FOR{each row $i = 1, \dots, N$ {\bfseries in parallel}}
    \STATE Sort $A_{i,:}$ by $\mathrm{freq}(\cdot)$ in descending order
  \ENDFOR
  \STATE {\bfseries Phase 2: Iterative Refinement}
  \FOR{$t = 1, \dots, T$}
    \STATE Scoreboard $\mathrm{cnt}(j,v) = |\{i : A_{i,j} = v\}|$ for each column $j$, value $v$
    \FOR{each value $v$}
      \STATE $\mathrm{pref}(v) = $ columns sorted by $\mathrm{cnt}(\cdot, v)$ descending
    \ENDFOR
    \FOR{each row $i = 1, \dots, N$ {\bfseries in parallel}}
      \STATE $Q \gets$ max-heap, $\mathrm{used} \gets \emptyset$
      \FOR{each item position $k = 1, \dots, m$}
        \STATE $v \gets A_{i,k}$, top-preference column $j^{\star} \gets \mathrm{pref}(v)[1]$
        \STATE Insert $(\mathrm{cnt}(j^{\star}, v),\; k,\; 1)$ into $Q$
      \ENDFOR
      \WHILE{$Q \neq \emptyset$}
        \STATE $(\mathit{score},\, k,\, p) \gets $ extract max from $Q$
        \STATE $v \gets A_{i,k}$, target column $j \gets \mathrm{pref}(v)[p]$
        \IF{$j \notin \mathrm{used}$}
          \STATE Assign $\hat{A}_{i,j} \gets v$ and add $j$ to $\mathrm{used}$
        \ELSE
          \STATE Insert $(\mathrm{cnt}(\mathrm{pref}(v)[p\!+\!1],\, v),\; k,\; p\!+\!1)$ into $Q$
        \ENDIF
      \ENDWHILE
      \STATE $A_{i,:} \gets \hat{A}_{i,:}$
    \ENDFOR
    \STATE $H^{(t)} \gets \sum_{j=1}^{m} H(A_{:,j})$
    \IF{$t > 1$ {\bfseries and} $|H^{(t)} - H^{(t-1)}| / H^{(t-1)} < \varepsilon$}
      \STATE Early terminate and return $A$
    \ENDIF
  \ENDFOR
\end{algorithmic}
\end{algorithm}



\section{Hardness}
\label{sec:hardness}

The decision version \textsc{Col-Entropy} asks, given $A$ and a threshold $T$,
whether some tuple of row permutations $\sigma$ achieves $F(\sigma) \geq T$
for $F$ as defined in Eq.~\eqref{eq:objective}.

\begin{theorem}\label{thm:np-complete}
\textsc{Col-Entropy} is NP-complete for every fixed $m \geq 3$.
\end{theorem}

\begin{proof}
Membership is immediate: $\sigma$ itself is a certificate verifiable in
$O(Nm)$ by recomputing $F(\sigma)$ from the counts $c_v^{(j)}$.

For hardness, we reduce from \textsc{Graph $m$-Coloring}~\cite{karp1972}.
Given $G = (V, E)$ with $N = |E|$, build $A_G \in \mathbb{Z}^{N \times m}$ as
follows: for each edge $(u, v) \in E$, the corresponding row contains $u$,
$v$, and $m - 2$ unique dummy items that appear nowhere else in $A_G$. Set
$T = \sum_{u \in V} \deg(u) \log \deg(u)$. We show that $F^\ast(A_G) = T$ iff
$G$ is $m$-colorable.

\emph{Only if.} Given a proper coloring $c : V \to [m]$, assign each vertex
within each row to the column matching its color, and distribute the dummies
arbitrarily among the $m - 2$ remaining columns. Since adjacent vertices
receive different colors, no column is used twice per row. Every vertex $u$
lands entirely in column $c(u)$ and contributes $\deg(u) \log \deg(u)$;
dummies have frequency one and contribute zero. Hence $F(\sigma) = T$.

\emph{If.} By strict convexity of $f(x) = x \log x$, any distribution of a
vertex $u$'s $\deg(u)$ copies across columns satisfies
$\sum_j f(c_u^{(j)}) \leq f(\deg(u))$, with equality only when all copies lie
in a single column. Summing over vertices and using $f(1) = 0$ for dummies
yields $F(\sigma) \leq T$ for every $\sigma$; in particular $F^\ast(A_G) \leq T$,
so any YES-instance ($F^\ast(A_G) \geq T$) must attain equality. Equality at
$T$ forces each vertex into a unique column $c(u)$; since the two endpoints
of every edge share a row and thus land in different columns, $c$ is a proper
$m$-coloring of $G$.
\end{proof}

Since $F^\ast(A_G) \leq T$ holds unconditionally, any algorithm returning
$F^\ast(A)$ decides $m$-colorability via a single comparison against $T$.
Thus the optimization version---maximizing $F$ (equivalently, minimizing
$\sum_j H(A_{:,j})$)---is also NP-hard for every fixed $m \geq 3$.\footnote{We note that 2-coloring, the decision problem of determining if a graph is bipartite, is in P, which is why we require $m \ge 3$ for our reduction. In other words, \textsc{Col-Entropy} is solvable in polynomial time when we have fewer than 3 columns.} 


\section{Experiments}

Present experiment results
\section{Amazon Product Search}

Amazon Search powers product discovery for hundreds of millions of customers, processing billions of search queries per week across 
global marketplaces. Each query triggers a complex retrieval and ranking pipeline that identifies relevant products from a catalog of 
hundreds of millions of items. With the increasing adoption of large language models for query understanding, result generation, and re
-ranking, the computational cost of serving queries in real time has grown substantially. Precomputing and caching results is therefore
a critical strategy for maintaining low-latency responses at scale. We evaluate CARAMEL on two result caching workloads derived from 
Amazon's US marketplace.

Query rewriting cache. A common component in search systems is query rewriting, where the system maps an incoming query to a set of 
alternative or expanded queries that are also issued to the retrieval engine. For example, a query "wireless earbuds" may be rewritten 
to {"bluetooth earbuds", "wireless headphones", "earbuds noise cancelling"}. These rewrite mappings are expensive to compute — often 
involving LLM-based models — and are therefore precomputed and stored in a query-to-query cache. We extract the query frequency 
distribution from a 5-week observation window of Amazon US mobile search traffic, filtered to keyword searches with bot and spam 
traffic removed. The dataset contains 792 million unique queries with a total volume of 3.3 billion occurrences. The distribution is 
heavily skewed: 79.6\% of queries appear exactly once, while the most popular query occurs 1.9 million times. In a query rewriting 
cache, popular rewrite targets (e.g., common product-type queries like "laptop" or "running shoes") appear in the value lists of many 
different source queries, creating the item-level repetition that CARAMEL is designed to exploit.

Query-to-ASIN result cache. The primary output of a search system is a ranked list of products (ASINs) for each query. Precomputing 
these result lists offline and serving them from a cache avoids the cost of real-time retrieval and ranking. We obtain the ASIN click 
frequency distribution from a one-week subsample of US production search traffic, covering 12 million unique ASINs that received at 
least one click, with approximately 150 million total clicks. The distribution follows a log-normal pattern: 48\% of ASINs receive 
exactly one click, and the top 1\% of ASINs account for 37\% of all clicks. This concentration means that popular products appear across 
the result lists of many different queries, providing significant redundancy for compression.
\section{Related Work}

Discuss prior work in result caching for recommender systems. 
\section{Conclusion}

Wrap it up...



\bibliographystyle{ACM-Reference-Format}
\bibliography{main}
\end{document}
\endinput
%%
%% End of file `sample-sigconf.tex'.
